\section{Introduction}
\label{sec:intro}

Make-on-demand chemical libraries now exceed \num{e10} compounds, far beyond
what structure-based docking can exhaustively evaluate. Ultra-large docking
campaigns have shown that screening deeper into such libraries recovers
better and more diverse chemotypes~\citep{lyu2019ultralarge}, but the cost of
docking every molecule grows with library size. The dominant response is to
dock a subset and train a \emph{surrogate} model that predicts the docking
score of the remainder from a cheap molecular representation, then dock only
the molecules the surrogate ranks highly. This is the basis of pool-based
active learning for virtual screening~\citep{graff2021molpal} and of related
iterative schemes~\citep{gentile2022deepdocking}.

The surrogate's representation is therefore a consequential design choice, and
it has become common to use a pretrained molecular foundation model for it.
Such models are available in increasing number and variety --- trained on
SMILES~\citep{ross2022molformer,soares2025smited}, on
SELFIES, on molecular graphs~\citep{rong2020grover}, on 3D
conformers~\citep{zhou2023unimol}, and on 2D depictions --- and recent work has
systematically benchmarked how combining them affects downstream prediction
tasks in drug discovery~\citep{martinezgalindo2026multimodality}. That study
trained over \num{1400} models across \num{29} benchmarks and concluded that
additional modalities generally help provided an effective fusion strategy is
used, with late fusion outperforming intermediate fusion.

Two gaps motivate the present work.

\paragraph{The missing baseline.}
Benchmarks of learned molecular representations are typically compared against
each other, or against a fingerprint, rather than against cheap
physicochemical descriptors. For docking-score prediction this omission is
consequential, because docking scoring functions are not scale-free: a
DOCK3.7 score is a sum of per-atom interaction terms and correlates strongly
with molecular size. A representation that encodes little beyond size can
therefore post a high correlation against the true score while carrying no
information about binding. Without a descriptor control, it is not possible to
tell these cases apart.

\paragraph{The missing size control.}
The same mechanism contaminates the evaluation metrics themselves. Pearson
correlation and RMSE computed over a whole pool are dominated by the
size-driven component of the variance. A screener, however, operates within a
desired physicochemical envelope and needs a model that discriminates
\emph{among} molecules of comparable size. Metrics that do not control for size
do not answer that question.

\paragraph{Contributions.}
\begin{enumerate}
  \item We establish a \num{12}-feature physicochemical descriptor baseline for
        docking-score regression, evaluated under identical splits, prediction
        head and metrics as the learned representations it is compared
        against (Section~\ref{sec:methods}).
  \item We introduce a size-controlled evaluation protocol --- partial
        correlation, residual correlation after a polynomial size fit, and
        metrics within equal-count heavy-atom strata --- and show it separates
        genuine chemical signal from size encoding
        (Section~\ref{sec:results-size}).
  \item We benchmark \todo{N} pretrained molecular embeddings, spanning
        \numrange{512}{3400} dimensions, against that baseline on pools of up to
        \num{116} million molecules, reporting learning curves rather than
        single operating points (Section~\ref{sec:results-backbones}).
  \item We examine late fusion across representations and show that the
        prediction-averaging strategy reported as effective in the multimodal
        benchmarking literature degrades sharply when member quality is uneven
        (Section~\ref{sec:results-fusion}).
  \item We release a fully reproducible pipeline, including the dataset
        alignment verification that this benchmark turns out to require
        (Section~\ref{sec:repro}).
\end{enumerate}

\todo{Finalise the headline framing once the backbone sweep completes. If
learned embeddings fail to clear the descriptor baseline, the paper is a
negative result about representation quality and the title should say so. If
they clear it comfortably, the contribution is the size-controlled protocol
plus a cost--accuracy analysis, and the framing shifts accordingly.}
